\abstractpagero{Extras}{%
  Prezenta lucrare de diplomă abordează problema recunoașterii automate a
  instrumentelor muzicale pe baza înregistrărilor audio, utilizând metode
  moderne de învățare automată. Sistemul dezvoltat este capabil să clasifice
  în timp real șapte categorii de instrumente (chitară, pian, voce,
  instrumente cu coarde, instrumente cu ancie, instrumente de alamă și
  zgomot ambiental) pornind de la semnale audio monofonice.

  Abordarea propusă se bazează pe extragerea a trei tipuri de caracteristici
  spectrale (Log-Mel spectrograme, STFT și MFCC) și pe antrenarea unor rețele
  neuronale convoluționale bidimensionale (2D~CNN) în mediul Google Colab.
  Setul de date hibrid a fost creat prin combinarea unor eșantioane curate
  descărcate de pe internet, a unor re-înregistrări proprii prin microfon și
  a zgomotului ambiental din baza de date ESC-50, iar un modul de augmentare
  a datelor simulează condițiile reale de captare prin adăugarea de reverberație,
  zgomot ambiental și distorsiuni de microfon.

  Rezultatele experimentale arată o acuratețe competitivă a modelului pe
  datele de test, iar sistemul a fost validat prin recunoașterea în timp real
  a sunetelor captate cu microfonul laptopului. Lucrarea contribuie la
  demonstrarea fezabilității unui clasificator robust de instrumente muzicale,
  antrenat cu resurse computaționale accesibile.
}{%
\textit{\textbf{Cuvinte cheie:}} recunoașterea instrumentelor muzicale, rețele
neuronale convoluționale, Mel-spectrogramă, STFT, MFCC, clasificare audio, formă de undă, imagine spectrală%
}
